\subsubsection{弱收敛}\label{subsec:weak-convergence}

求解微分方程的过程通常涉及一列近似问题的解. 往往可以在某个 Banach 空间的范数下得到这些近似解的一致界, 但要证明它们确实在该 Banach 空间中收敛并不那么容易. 在这些情况下, 一种称为弱收敛的较弱收敛概念极为有用. 函数序列的弱收敛意味着, 某些线性泛函作用于该序列时, 得到收敛的数列. 当定义分布序列 $f_n$ 在对每个试验函数 $\phi$ 都有 $(f_n,\phi)$ 收敛时为收敛序列时, 已经遇到过这样的概念. 在 Banach 空间中, 作如下定义.

\begin{Definition}
\label{def:weak-and-weak-star-convergence}
设 $X$ 为 Banach 空间. 若对每个 $f\in X^*$, $f(x_n)$ 都收敛到 $f(x)$, 则称 $X$ 中的序列 $x_n$ \textbf{弱收敛}到 $x$. 若对每个 $x\in X$, $f_n(x)$ 都收敛到 $f(x)$, 则称 $X^*$ 中的序列 $f_n$ \textbf{弱-* 收敛}到 $f$.
\end{Definition}

为区分记号, 将范数收敛写为 $x_n\to x$, 弱收敛写为 $x_n\rightharpoonup x$, 弱-* 收敛写为 $x_n\overset{*}{\rightharpoonup}x$.

\begin{Remark}
\label{rem:weak-convergence-consequences}
这里列出弱收敛定义的若干重要推论. 对这些断言的证明 (通常很短), 只给出极简略的提示.\begin{itemize}
\item 由定义直接得到, $X^*$ 中的弱收敛序列也弱-* 收敛. 反过来一般不成立, 但若 $X$ 自反, 则成立.
\item 弱极限的唯一性由 Hahn-Banach 定理得到; 弱-* 极限的唯一性由定义显然可知.
\item 由一致有界原理得到, 弱收敛或弱-* 收敛序列有界.
\item 若对每个 $f\in X^*$, $f(x_n)$ 都为 Cauchy 序列, 则
\[l(f):=\lim_{n\to\infty}f(x_n)\]
定义了 $X^*$ 上的一个线性泛函, 但未必知道该线性泛函可以表示为 $l(f)=f(x)$ 的形式, 其中 $x\in X$. 因而弱收敛未必满足 Cauchy 准则. 满足这一准则的空间称为弱 (序列) 完备的. 显然, 自反空间弱完备. 另一方面, 弱-* 收敛总满足 Cauchy 准则.
\end{itemize}
\end{Remark}

\Needspace{4\baselineskip}
\begin{Example}
\label{ex:orthonormal-basis-weak-zero-example}
设 $H$ 为可分无穷维 Hilbert 空间, $e_n$ 为 $H$ 的一组标准正交基. 那么 $e_n$ 弱收敛到零 (Problem~\ref{ex:orthonormal-basis-weak-zero}).
\end{Example}

\Needspace{4\baselineskip}
\begin{Example}
\label{ex:oscillatory-sines-weak-zero}
前一例的一个特殊情形是 $L^2(0,1)$ 中的序列 $f_n(x)=\sin nx$. 这一由振荡越来越剧烈的函数组成的序列弱收敛到零. 当然, 它并不以任何逐点意义收敛到零, 尽管可以说它 ``平均而言'' 收敛到零.
\end{Example}
